% ================================================================================
% PECA 130 / PAPEL DE TRABALHO PT-XX - ANALISE DE OUTLIERS DE RENDA NO CAF
% ================================================================================
% Data: 02/dezembro/2025
% Objetivo: Documentar analise estatistica de valores atipicos (outliers) de renda
%           declarada no CAF, utilizando metodologia Tukey e agregacao por UF.
% Referencia: peca130_pt_outliers_renda.pdf
% ================================================================================

\documentclass[11pt,a4paper]{article}

% ========== PACOTES PADRAO ==========
\usepackage[utf8]{inputenc}
\usepackage[brazil]{babel}
\usepackage[T1]{fontenc}
\usepackage{lmodern}
\usepackage[margin=2.5cm]{geometry}
\usepackage{graphicx}
\usepackage{float}
\usepackage{longtable}
\usepackage{booktabs}
\usepackage{xcolor}
\usepackage{colortbl}
\usepackage{fancyhdr}
\usepackage{amsmath, amssymb}
\usepackage{listings}
\usepackage{enumitem}
\usepackage[hidelinks]{hyperref}
\usepackage[breakable]{tcolorbox}

% ========== CORES TCU OFICIAIS ==========
\definecolor{tcured}{RGB}{204,0,0}      % Vermelho TCU #CC0000
\definecolor{tcublue}{RGB}{0,51,102}    % Azul TCU #003366
\definecolor{tcugreen}{RGB}{0,153,76}   % Verde TCU #006633

% ========== CONFIGURACAO DE CABECALHO ==========
\pagestyle{fancy}
\fancyhead[L]{Peca 130 / PT - Outliers de Renda no CAF}
\fancyhead[R]{TCU}
\setlength{\headheight}{14pt}

% ========== CONFIGURACAO DE LISTINGS ==========
\lstset{
    basicstyle=\ttfamily\small,
    frame=single,
    breaklines=true,
    numbers=left,
    numberstyle=\tiny\color{gray},
    keywordstyle=\color{tcublue}\bfseries,
    commentstyle=\color{tcugreen}\itshape,
    stringstyle=\color{tcured},
    columns=fullflexible
}

% ========== INFORMACOES DO DOCUMENTO ==========
\title{Papel de Trabalho de Auditoria Tecnica\\
\textbf{Peca 130: Analise Estatistica de Outliers de Renda no CAF}}
\author{Equipe de Auditoria --- Fiscalizacao CAF-MAPA}
\date{02 de dezembro de 2025}

\begin{document}

\maketitle

% ================================================================================
\section{Objetivo}
% ================================================================================

Este papel de trabalho documenta a \textbf{analise estatistica de valores atipicos (outliers)
de renda declarada} no Cadastro da Agricultura Familiar (CAF), com o objetivo de:

\begin{itemize}
    \item Identificar unidades familiares com rendas declaradas estatisticamente anomalas;
    \item Quantificar a extensao do problema de qualidade de dados relacionado a renda;
    \item Fornecer base documental para as constatacoes do relatorio de auditoria.
\end{itemize}

\textbf{Contexto}: A analise de outliers e fundamental para avaliar a confiabilidade dos
dados de renda no CAF, que sao utilizados para definicao de politicas publicas voltadas
a agricultura familiar.

% ================================================================================
\section{Metodologia}
% ================================================================================

\subsection{Base de Dados}

\begin{itemize}
    \item \textbf{Fonte}: Tabela \texttt{caf\_mapa.S\_RENDA}
    \item \textbf{Populacao}: 12.849.431 registros de renda
    \item \textbf{Unidades Familiares ativas}: 3.304.174 UFs
    \item \textbf{Filtros aplicados}:
    \begin{itemize}
        \item Apenas UFs com situacao ativa (\texttt{id\_tipo\_situacao\_unidade\_familiar = 1})
        \item \textbf{Exclusao}: Tipo de renda ``Autoconsumo UFPA'' (1.795.891 registros)
    \end{itemize}
\end{itemize}

\subsection{Justificativa da Exclusao do Autoconsumo}

O tipo de renda ``Autoconsumo UFPA'' representa o valor estimado da producao consumida
pela propria familia, \textbf{nao constituindo renda monetaria efetiva}. Sua inclusao
distorceria a analise de outliers, pois:

\begin{enumerate}
    \item Representa producao para subsistencia, nao transacao comercial;
    \item Os valores sao estimativas, nao declaracoes de receita;
    \item Soma-lo a renda monetaria criaria valores inflacionados artificialmente.
\end{enumerate}

\subsection{Unidade de Analise}

A analise foi realizada por \textbf{Unidade Familiar (UF)}, somando todas as rendas
declaradas de cada UF. Esta abordagem evita a duplicacao que ocorreria ao analisar
por individuo, quando multiplos membros da mesma familia reportam a mesma renda.

\subsection{Metodo de Deteccao de Outliers}

Foi utilizada a \textbf{Regra de Tukey} (fence method), metodo estatistico robusto
amplamente aceito:

\begin{align*}
\text{IQR} &= Q3 - Q1 \\
\text{Limiar Superior} &= Q3 + 1,5 \times \text{IQR}
\end{align*}

Onde:
\begin{itemize}
    \item $Q1$ = Primeiro quartil (25\% dos dados)
    \item $Q3$ = Terceiro quartil (75\% dos dados)
    \item IQR = Amplitude interquartil
\end{itemize}

Valores de renda acima do limiar superior sao classificados como \textbf{outliers moderados}.

\subsection{Query SQL Principal}

\begin{lstlisting}[language=SQL, caption={Calculo de estatisticas de renda por UF}]
WITH renda_uf AS (
    SELECT
        r.id_unidade_familiar,
        SUM(r.vl_renda_auferida) AS renda_total_uf
    FROM caf_mapa."S_RENDA" r
    INNER JOIN caf_mapa."S_UNIDADE_FAMILIAR" uf
        ON r.id_unidade_familiar = uf.id_unidade_familiar
    INNER JOIN caf_mapa."S_TIPO_RENDA" tr
        ON r.id_tipo_renda = tr.id_tipo_renda
    WHERE uf.id_tipo_situacao_unidade_familiar = 1
      AND tr.ds_tipo_renda != 'Autoconsumo UFPA'
    GROUP BY r.id_unidade_familiar
)
SELECT
    ROUND(PERCENTILE_CONT(0.25) WITHIN GROUP
        (ORDER BY renda_total_uf)::NUMERIC, 2) AS q1,
    ROUND(PERCENTILE_CONT(0.50) WITHIN GROUP
        (ORDER BY renda_total_uf)::NUMERIC, 2) AS mediana,
    ROUND(PERCENTILE_CONT(0.75) WITHIN GROUP
        (ORDER BY renda_total_uf)::NUMERIC, 2) AS q3,
    ROUND(AVG(renda_total_uf)::NUMERIC, 2) AS media,
    COUNT(*) AS total_ufs
FROM renda_uf;
\end{lstlisting}

% ================================================================================
\section{Resultados}
% ================================================================================

\subsection{Estatisticas Descritivas}

\begin{table}[H]
\centering
\caption{Estatisticas de Renda Total por Unidade Familiar}
\label{tab:estatisticas}
\begin{tabular}{lrl}
\toprule
\textbf{Parametro} & \textbf{Valor} & \textbf{Observacao} \\
\midrule
Primeiro Quartil (Q1) & R\$ 9.440,00 & 25\% das UFs \\
\textbf{Mediana (Q2)} & \textbf{R\$ 18.830,75} & Valor central \\
Terceiro Quartil (Q3) & R\$ 48.874,75 & 75\% das UFs \\
\midrule
Amplitude Interquartil (IQR) & R\$ 39.434,75 & Q3 - Q1 \\
\textbf{Limiar Tukey} & \textbf{R\$ 108.026,88} & Q3 + 1,5 $\times$ IQR \\
\midrule
Media & R\$ 58.307,00 & Influenciada por outliers \\
Total de UFs analisadas & 3.304.174 & 100\% \\
\bottomrule
\end{tabular}
\end{table}

\subsection{Outliers Identificados}

\begin{table}[H]
\centering
\caption{Classificacao de Outliers}
\label{tab:outliers}
\begin{tabular}{lrr}
\toprule
\textbf{Categoria} & \textbf{Quantidade} & \textbf{Percentual} \\
\midrule
Renda normal ($\leq$ R\$ 108.027) & 2.809.875 & 85,04\% \\
\midrule
\textbf{Outliers Tukey} (R\$ 108k - 200k) & 231.883 & 7,02\% \\
Outliers altos (R\$ 200k - 500k) & 254.317 & 7,70\% \\
Outliers muito altos (R\$ 500k - 1M) & 7.192 & 0,22\% \\
\textbf{Outliers extremos} (R\$ 1M - 10M) & 768 & 0,02\% \\
\textbf{Outliers absurdos} ($>$ R\$ 10M) & 139 & $<$0,01\% \\
\midrule
\textbf{Total de outliers Tukey} & \textbf{494.299} & \textbf{14,96\%} \\
\textbf{Total $>$ R\$ 1 milhao} & \textbf{907} & 0,03\% \\
\bottomrule
\end{tabular}
\end{table}

\subsection{TOP 10 Rendas Mais Extremas}

\begin{table}[H]
\centering
\caption{Maiores Outliers de Renda Declarada}
\label{tab:top10}
\begin{tabular}{crr}
\toprule
\textbf{Rank} & \textbf{Renda Anual Declarada} & \textbf{Renda Mensal Implicita} \\
\midrule
1 & R\$ 167.320.132,00 & R\$ 13.943.344,33 \\
2 & R\$ 105.347.730,00 & R\$ 8.778.977,50 \\
3 & R\$ 97.322.000,00 & R\$ 8.110.166,67 \\
4 & R\$ 92.228.407,00 & R\$ 7.685.700,58 \\
5 & R\$ 78.912.359,00 & R\$ 6.576.029,92 \\
6 & R\$ 78.258.000,00 & R\$ 6.521.500,00 \\
7 & R\$ 68.570.000,00 & R\$ 5.714.166,67 \\
8 & R\$ 67.921.000,00 & R\$ 5.660.083,33 \\
9 & R\$ 63.825.000,00 & R\$ 5.318.750,00 \\
10 & R\$ 62.025.822,00 & R\$ 5.168.818,50 \\
\bottomrule
\end{tabular}
\end{table}

\textbf{Observacao}: O maior outlier representa uma renda mensal implicita de
\textbf{R\$ 13,9 milhoes}, valor incompativel com o perfil da agricultura familiar.

% ================================================================================
\section{Analise}
% ================================================================================

\subsection{Significancia Estatistica}

O percentual de \textbf{14,96\% de outliers} e estatisticamente significativo e indica
problemas sistemicos na validacao de dados de renda no CAF:

\begin{itemize}
    \item Pela regra de Tukey, espera-se que outliers representem uma fracao pequena
    (tipicamente $<$5\%) em distribuicoes normais;
    \item A presenca de quase 15\% de outliers sugere: (a) erros de digitacao,
    (b) falta de validacao de limites, ou (c) declaracoes fraudulentas.
\end{itemize}

\subsection{Implicacoes para Politicas Publicas}

A agricultura familiar e definida, entre outros criterios, por limites de renda.
Outliers extremos como rendas acima de R\$ 1 milhao (907 casos) sao \textbf{incompativeis}
com o conceito de agricultura familiar e podem indicar:

\begin{enumerate}
    \item \textbf{Erros de digitacao}: Valores inseridos sem ponto decimal (ex: 18900 ao
    inves de 189,00);
    \item \textbf{Falta de validacao}: O sistema nao rejeita valores implausveis;
    \item \textbf{Declaracoes incorretas}: Informacoes nao condizentes com a realidade.
\end{enumerate}

\subsection{Comparacao com Limite de Enquadramento}

O limite de renda para enquadramento como agricultor familiar (Lei 11.326/2006) e de
aproximadamente R\$ 500.000 anuais. Identificamos:

\begin{itemize}
    \item \textbf{8.099 UFs} com renda acima de R\$ 500.000 (0,25\% do total)
    \item \textbf{907 UFs} com renda acima de R\$ 1.000.000 (0,03\% do total)
\end{itemize}

Estes casos requerem verificacao para confirmar enquadramento legal.

% ================================================================================
\section{Conclusao}
% ================================================================================

A analise estatistica de outliers de renda no CAF revela \textbf{problemas significativos
de qualidade de dados}:

\begin{enumerate}
    \item \textbf{494.299 unidades familiares} (14,96\%) apresentam rendas consideradas
    outliers pela regra de Tukey;

    \item \textbf{907 UFs} declararam rendas acima de R\$ 1 milhao, com o caso mais
    extremo atingindo R\$ 167 milhoes --- valor claramente incompativel com agricultura
    familiar;

    \item A \textbf{ausencia de validacao de limites} no sistema permite a insercao de
    valores implausveis, comprometendo a confiabilidade estatistica da base;

    \item Os outliers \textbf{distorcem a media} de renda: com todos os registros,
    a media e R\$ 58.307; excluindo apenas os 139 casos acima de R\$ 10 milhoes,
    a media cai para R\$ 57.129 (distorcao de +2,06\%). A diferenca expressiva
    entre media (R\$ 58.307) e mediana (R\$ 18.831) evidencia a assimetria causada
    pelos outliers.
\end{enumerate}

\textbf{Recomendacao}: Implementar controles preventivos de validacao de renda no CAF,
incluindo alertas para valores acima de limiares razoaveis e verificacao cruzada com
bases externas (ex: Receita Federal).

% ================================================================================
\section{Anexos}
% ================================================================================

\subsection{Parametros de Conexao}

\begin{verbatim}
Host: 172.23.80.1
Port: 5433
Database: caf_mapa
User: postgres
Data de execucao: 02/dezembro/2025
\end{verbatim}

\subsection{Query de Outliers Tukey}

\begin{lstlisting}[language=SQL, caption={Identificacao de outliers pela regra de Tukey}]
WITH renda_uf AS (
    SELECT
        r.id_unidade_familiar,
        SUM(r.vl_renda_auferida) AS renda_total_uf
    FROM caf_mapa."S_RENDA" r
    INNER JOIN caf_mapa."S_UNIDADE_FAMILIAR" uf
        ON r.id_unidade_familiar = uf.id_unidade_familiar
    INNER JOIN caf_mapa."S_TIPO_RENDA" tr
        ON r.id_tipo_renda = tr.id_tipo_renda
    WHERE uf.id_tipo_situacao_unidade_familiar = 1
      AND tr.ds_tipo_renda != 'Autoconsumo UFPA'
    GROUP BY r.id_unidade_familiar
),
quartis AS (
    SELECT
        PERCENTILE_CONT(0.25) WITHIN GROUP
            (ORDER BY renda_total_uf) AS q1,
        PERCENTILE_CONT(0.75) WITHIN GROUP
            (ORDER BY renda_total_uf) AS q3
    FROM renda_uf
)
SELECT
    COUNT(*) AS outliers_tukey,
    ROUND(COUNT(*)::NUMERIC * 100 /
        (SELECT COUNT(*) FROM renda_uf), 3) AS percentual
FROM renda_uf, quartis
WHERE renda_total_uf > q3 + 1.5 * (q3 - q1);
-- Resultado: 494.300 outliers (14,96%)
\end{lstlisting}

\subsection{Tipos de Renda Analisados}

Foram considerados 25 tipos de renda (excluindo Autoconsumo UFPA):

\begin{itemize}[noitemsep]
    \item Lavouras Temporarias (3.811.686 registros)
    \item Rendas fora do estabelecimento (2.773.381)
    \item Producao Animal (2.315.362)
    \item Lavouras Permanentes (930.697)
    \item E outros 21 tipos menores
\end{itemize}

\subsection{Artefatos Gerados}

\begin{itemize}
    \item \texttt{peca130\_pt\_outliers\_renda.tex} --- Este documento (codigo fonte)
    \item \texttt{peca130\_pt\_outliers\_renda.pdf} --- Versao compilada (Peca 130)
\end{itemize}

\end{document}
